% Source: https://artofproblemsolving.com/community/c3h2496532_number_theory
% Stable identity: aops-topic-2496532-q01
The three-digit multiples of $11$ are $11\cdot10,\dots,11\cdot90$, which is $81$ numbers. For $n=\overline{abc}$, $11\mid n$ exactly when $11\mid a-b+c$; since $-8\le a-b+c\le18$, this means $b=a+c$ or $b=a+c-11$. Now count the multiples with a repeated digit.
\begin{itemize}
\item $a=b$: then $c=0$ (or $c=11$, impossible), giving $110,220,\dots,990$: $9$ numbers.
\item $b=c$: then $a=0$ or $a=11$, impossible.
\item $a=c$: then $b=2a$, giving $121,242,363,484$, or $b=2a-11$, giving $616,737,858,979$: $8$ numbers.
\end{itemize}
These lists do not overlap, so $17$ multiples have a repeated digit and $81-17=64$ have three distinct digits. Since $64$ is not among the first three choices, the answer is $\textbf{(D)}$ \emph{none of these} (the number is $\boxed{64}$).
